\section{Conclusion and Future Work}
\label{sec:conclusions}
% EDIT (2026-10-29): rewritten for the new structure; repeated numbers reduced to the ones that
% carry the argument.
We studied prompt design for LLM-driven GPU code optimization inside a closed loop that profiles,
generates, builds, verifies and measures, and we pushed that loop past single-kernel benchmarks onto
two production scientific applications.

Three results stand out. First, guardrails decide what the models do with the freedom they are given.
Which instructions the prompt carries decides how often a model returns a valid candidate at the first
attempt and how slow its worst result is, and an explicit semantic-preservation constraint is what
separates an honest result from a gamed one: without it, both models delete the computation a
benchmark exists to measure and report speedups of three orders of magnitude that survive the
benchmark's own correctness check. Second, profiling data earns its place in the prompt in proportion
to how hard the bottleneck is to see from the source. On single-kernel benchmarks, where the hot loop
is obvious, it moves individual results in both directions with no net gain; on a 21k-line application
where it is not, it is worth $7.7$--$11.5\%$ to both models. Interpreting that data into named
remedies does not improve the average; it widens the spread of outcomes. Third, and most consequential
for anyone building such a system, the models reliably find local, exact transformations and reliably
miss cross-cutting restructuring: on Cholla every model output under the baseline prompt beat our
hand-tuned reference, while on AutoDock-GPU every candidate left in place the seven block reductions
whose fusion the hand reference exploits.

The failures were as informative as the speedups. Editing a single file inside a real program means the
model cannot see how that file is compiled, and three of our application failures follow directly.
Applications also swallow the diagnostics a repair loop needs, and shared machines silently corrupt
measurements unless contention is detected per repetition. None of these appear when optimizing
self-contained benchmark kernels, which is precisely why benchmark-only evaluation overstates how ready
such systems are.

Two directions follow. The prompt should carry the file's \emph{compilation context}---whether it is a
standalone translation unit, which symbols are already in scope, and which declarations other files
depend on---since that is cheap to extract and addresses the dominant real-code failure mode. And
profiling data should be supplied as \emph{evidence rather than remedy}: the measured condition,
not the transformation it suggests, so that the model keeps the choice that the checklist took from
it.
